\paragraph{Construction of the main summands}
In rank \(r\), take a nonzero good form \(z\) with the given Hecke eigensystem on \(W_m\).
Since \(e_0z=z\) and \(e_0\) is a sum of the idempotents of the admissible label tuples,
one of those tuples acts nontrivially on \(z\). The isomorphism to algebraic forms intertwines
right translation, hence also the Hecke action away from the assigned level.
At the assigned places, principal level makes the maximal compact group act through the general
linear group of the residue field. The nonzero tuple action shows that each constituent label
occurs. Attach \(R\) to this algebraic eigenform by the standard local--global compatibility
theorem. At the assigned places the label roots describe the inertia action, and at their
conjugates the local conditions follow by conjugate duality. At the other places spherical
compatibility gives unramifiedness, or crystallinity and the specified Hodge--Tate jumps at \(\ell\).
The first Hecke eigenvalue and the trace of the attached representation are obtained by negating
the next-to-leading coefficients. Local--global compatibility therefore gives
\(\operatorname{tr}R(g)\) as the limit of the transported first Hecke eigenvalues as Frobenius
elements tend to \(g\). The systems needed on every \(W_j\) occur on \(W_m\), by the comparison
of good eigenspaces, which gives vanishing on \(W_j\) when the eigenspace on \(W_m\) is zero.
The lifts of the good eigencomponents span the image of each lifting map.
\leanlinks{\href{\upstream/grh/src/New/Congruences/GaloisRepresentations/RowsFromEigensystems.lean\#L480}{Constructing the attached representation of an eigensystem}
  \leanlinksep \href{\upstream/grh/src/New/Congruences/GaloisRepresentations/AttachedRepresentations.lean\#L641}{The labels and local conditions}
  \leanlinksep \href{\upstream/grh/src/New/Congruences/FitTwistsExist.lean\#L326}{The eigensystems required for the lifts}}

Write \(r'\) for the source rank, \(t'\) for the rank difference, and \(\mu\) for the
character of the untwisted descent pair, of weight \(w\). Both ranks \(r',t'\) are odd, so
\(A=(r'-t')/2\) is an integer. The main representation \(R_x\) is \(R\) times the avatar of
\(x\mapsto\mu^{-1}(x/\bar x)|x|_E^A\).
A jump \(u\) at an embedding \(\sigma\) above \(v\) thus becomes
\(u-(A-\varepsilon_\sigma w(v))\), where \(\varepsilon_\sigma=1\) on the CM type
and \(-1\) on its conjugate. The tame part at the finite places scales the label roots by
the finite-field element naming the inverse tame character.
The cyclotomic power shifts the duality exponent to \(1-t'\).

In rank \(r\), expanding the \(q\)-th diagram weight gives the shifted jump
\(\lfloor t/2\rfloor-\varepsilon_\sigma(\lambda_q+w_j)\), \(0\le q<r\).
Here \(w_j\) is minus the weight of the character of stage \(j\), and \(\lambda_q\)
is the corresponding source parameter. Write \(G\) for the gap in the weights.
For \(q+1<k\), \(|\lambda_q|\le k/2\) gives \(|\varepsilon_\sigma u+w_j|\le t\).
For \(k\le q+1\le r\), put \(i=q+2-k\). Then
\(\lambda_q+w_j=(j-i+\tfrac12)G\), giving \(G/4\le |u|\le 2mG\).
There are \(m\) of the latter jumps, within the bound \(rm/k\) for the number of small jumps.
The two jumps at conjugate embeddings for the same \(q\) sum to \(2\lfloor t/2\rfloor\),
so at most \(r\) jumps at the pair are negative.
For rank-one main pairs the starting \(R\) is trivial, with jump \(0\), label \(X-1\),
and first eigenvalue \(1\). A main pair \((\nu,u_0)\) of weight \(-c\) has
\(t'=2u_0-1\) and shifted jump \(u_0-1-\varepsilon_\sigma c\).
The range \(m<u_0\le r\), \(w_j+r-t\le c\le w_j\) gives
\(|\varepsilon_\sigma(u_0-1-\varepsilon_\sigma c)+w_j|\le t\).
There are no small jumps, and the sum at the conjugate pair is \(2(u_0-1)\).
\leanlinks{\href{\upstream/grh/src/New/Congruences/GaloisRepresentations/RowsFromEigensystems.lean\#L143}{Bounds for the jumps in ranks \(r\) and one}
  \leanlinksep \href{\upstream/grh/src/New/Congruences/FitTwistsExist.lean\#L265}{The trivial starting representation in rank one}}

To obtain the full integral matrix algebra, use a level prime different from \(\ell\).
At one assigned place above it the scaled label is primitive in rank \(r\).
In rank one use the scaled \(X-1\) label.
The eigenvalues at a tame inertia element \(s\) are images of label roots, and a Frobenius
\(\phi\) conjugates \(R_x(s)\) to its \(q\)-th power, \(q=\mathrm Nw\).
In rank \(r=3^{d+1}\) we have \(q\equiv1\pmod9\).
Writing \(q^{3^d}=1+9h\) gives \((q^r-1)/(q^{3^d}-1)=81h^2+27h+3\).
The quotient is divisible by \(3\) but not by \(9\),
and exceeds \(3\), so it divides no power of a prime. Hence an \(\ell\)-power of a primitive
\((q^r-1)\)-st root has Frobenius period \(r\), since its period cannot divide \(3^d\).
The resulting \(\ell\)-adic eigenvalues have unit differences. Over the residue field,
apply powers of the inverse Frobenius to a single eigenvector to obtain a basis of eigenvectors.
Applying a linear relation among the matrices of \(\phi^{-a}s^b\), \(0\le a,b<r'\),
to each basis vector forces, for each \(a\), a polynomial of degree less than \(r'\) to
vanish at all \(r'\) distinct eigenvalues. These \(r'^2\) matrices are independent on reduction,
so their determinant is a unit and they are a basis of the integral matrix algebra.
Moreover \(\phi^{r'}\) commutes with the images of \(s,\phi\), hence with this basis, and acts
by a scalar. The product of the roots of the primitive label has order \(q-1\) and is a common
power of all the roots. In rank one the scaled root has the order of a power of the fixed tame
character. Thus a power of \(s\) acts by a scalar, and the prescribed tame order makes its
\(\ell\)-power multiples of the exponents \(1,\ldots,4r'+2\) remain separated from \(1\) on reduction.
For the determinant estimates at the other level places and their conjugates, the roots of the
labels have bounded orders, giving the bounds on the valuations of \(\det(R_x(s)-I)\).
\leanlinks{\href{\upstream/grh/src/New/Congruences/GaloisRepresentations/RowsAndColumnsFromLabels.lean\#L356}{Common powers and separation of label roots}
  \leanlinksep \href{\upstream/grh/src/New/Congruences/GaloisRepresentations/RowsAndColumnsFromLabels.lean\#L510}{The integral span from tame inertia and Frobenius}}

Finally take the reciprocal Hecke polynomial of a lift twisted by \(\vartheta\).
It is the product of \(\prod_{a<t'}(1-q^a\alpha X)\) and the source reciprocal polynomial
evaluated at \(\beta X\). If \(\alpha,\alpha'\) are the normalized uniformizer values of the
characters on the source and target, then \(\beta=q^{t'}\alpha/\alpha'\).
The constant terms are \(1\), so the negative linear coefficient is
\(\beta\,p(w,1)+\alpha\sum_{a<t'}q^a\).
Normalize by \(\alpha^{-1}\) and use local Artin reciprocity to transport the character values.
The first term is the avatar of the inverse twist times the avatar of the algebraic character above
times \(\iota(p(w,1))\). The second consists of the values of \(\epsilon_{\mathrm{cyc}}^{-a}\)
in the Frobenius convention used here. Local--global compatibility, continuity of the avatars,
and Chebotarev approximation of every \(g\) by Frobenius elements give the trace formula above.
\leanlinks{\href{\upstream/grh/src/New/Congruences/GaloisRepresentations/RowsFromEigensystems.lean\#L295}{The linear coefficient and the trace at Frobenius elements}
  \leanlinksep \href{\upstream/grh/src/New/Congruences/GaloisRepresentations/RowsFromEigensystems.lean\#L391}{Passing from Frobenius elements to every Galois element}}
